\cotefile{64}
\scanpage{40}
\mspage{\unsure{I c}}
% archive foliation: 39

\margin{\uline{NB} $X|\underline J\simeq\Delta$ est la configuration octaédrale\ldots}

\diagram{\begin{tikzcd}[ampersand replacement=\&,column sep=large,row sep=large]
\begin{gathered}(X,\underline I)\quad
\text{\deleted{\unsure{étale}}}\ \inserted{\underline I}\ \text{rev. étale fini de degré }4\end{gathered}
\arrow[d,leftrightarrow] \& \\
\begin{gathered}
\deleted{(\underline{\unsure{J}}(\Sigma_{\underline J}),\ t\in\Gamma(\Sigma^*_{\underline J}/S),\ \underline I\ \text{\deleted{\unsure{rev.}} torseur sous})}\\
(\underline I,\ \text{d'où }\underline J,\ t\in\Gamma(\Sigma^*_{\underline J}/S))
\end{gathered}\arrow[d,leftrightarrow]\arrow[r,leftrightarrow]
\&\begin{gathered}
\text{\uline{NB}}\quad\underline t\text{ et }\underline I\text{ indépendants}\\
(\underline J\ (\text{d'où }\underline V_{\underline J}\text{ et }\Sigma_J),\
\text{section }\inserted{t}\text{ de }\Sigma^*_J\text{ sur }S,\\
\text{torseur }I\text{ sur }\underline V_{\underline J})
\end{gathered}\\
\begin{gathered}
(\underline J,\ t\in\Gamma(\Sigma^*_{\underline J}/S),\
X\text{ rev. de }\Sigma^{\unsure{*}}_J,\ \deleted{\unsure{p}}\
\text{étale sur }\Sigma^*_J,\ \text{propre lisse sur }S,\\
\text{en chaque fibre géom. rev. \uline{abélien}}
\ \deleted{\text{\unsure{de degré } }4})\\
\inserted{\text{\unsure{fin. ram.}, \unsure{ramifiés, d’ordre} }2}\\
\text{\uline{NB}}\quad X\text{ et }t\text{ sont indép. l'un de l'autre}
\end{gathered}\arrow[d,leftrightarrow]\&\\
\begin{gathered}
(\underline J,\ t\in\Gamma(\Sigma^*_{\underline J}/S),\
\deleted{\unsure{I}}\underline K\text{ torseur sous }\underline J)\\
\text{\uline{NB}}\quad\underline K\text{ et }t\text{ sont indép. l'un de l'autre.}
\end{gathered}\&
\end{tikzcd}}

\margin{\deleted{$\unsure{\varphi_i\ \varphi_j\ \varphi_k}$}\quad
$\unsure{x_i}\quad x_j\quad x_k$
\diagram{\begin{tikzcd}[ampersand replacement=\&]
V_{\underline J}(M_0)\arrow[r]\&\unsure{\mathbb M_2}(M_0)\arrow[r,phantom,"\simeq"]\&\underline{\operatorname{End}}(M_0)\\
\&\&1\arrow[u,phantom,"\in" description]
\end{tikzcd}}}

\medskip\hrule\medskip
Quand $\underline I$ est muni d'une section, de sorte que $\underline I\simeq\underline J\amalg S$, alors on a un iso. can. $X\simeq\Sigma_J$, d'où $f_{\deleted{\unsure{I}}}:\Sigma_J\mathbin{\unsure{-|\longrightarrow}}\Sigma_J$.
\scanpage{41}
% archive foliation: 40
\begin{center}
UNIVERSITE DES SCIENCES\\
ET\\
TECHNIQUES DU LANGUEDOC
\end{center}
\hrule

Montpellier, le 9 mars 1982.

Service du Personnel Enseignant

MCS/FR.

N°1317.

Le Président

à

Mesdames et Messieurs\\
les enseignants de l'U.S.T.L.

\uline{OBJET} : Elections au Comité Technique Paritaire Central des Personnels Enseignants de statut universitaire -\\
Scrutin par correspondance de mars 1982.

\uline{Réf.} : Décret n°82.139 du 5 février 1982.\\
Arrêté du 8 février 1982.

J'ai l'honneur de vous adresser le matériel électoral nécessaire aux opérations de vote par correspondance, comprenant :

- 1 notice explicative ;

- 2 bulletins de vote (constitués par les listes des candidats)\\
\hspace*{2em}. SGEN - CFDT\\
\hspace*{2em}. SNESUP

- 2 professions de foi (1 pour chaque liste)

- jeu de 3 enveloppes

\hspace*{2em}. enveloppe n°1 blanche, ne comportant aucune marque ou distinction permettant d'en déterminer l'origine ;

\hspace*{2em}. enveloppe n°2 portant les mentions : nom, prénom, grade, affectation, signature ;

\hspace*{2em}. enveloppe n°3 libellée au nom de M. le Recteur, Président de la Commission de Dépouillement.

Je vous précise que les votes devront être adressés, \uline{par voie postale à M. le Recteur de l'Académie de Montpellier, la date limite de réception étant fixée au 30 mars 1982 à 17 heures,le cachet de la poste faisant foi}.

Les votes parvenus hors délais seront renvoyés aux électeurs.

Le dépouillement public se déroulera au Rectorat, Salle des Commissions, le \uline{31 mars 1982 à 14 heures}.

Le Président,

\unsure{L. Thaler}

\uline{L. THALER}.
\scanpage{42}
% archive foliation: 41
\uline{Prop} $V$ vectoriel de \unsure{rg} $2$ sur $\mathbb F_2$, donc $V\simeq V_J(\mathbb F_2)$, où $J=V^*$ ens. à trois éléments
\[
\underbrace{\operatorname{Aut}(V)}_{\Gamma}\xrightarrow{\sim}\unsure{\mathfrak S_J|\mathcal S_J}.
\qquad\text{Posons}\quad\unsure{\operatorname{Aut}(V)|\operatorname{Aut}^-(V)}\simeq\unsure{G_J^-|\mathfrak S_J^-}.
\]
\deleted{\unsure{Alors d’où|Alors on a}}
\[
\operatorname{Aut}(V)\subset\operatorname{End}(V)=M\deleted{\text{\unsure{can.}}}\quad(\text{alg. Azumaya}).
\]
Alors :

1) $\operatorname{Aut}^-(V)=V'{}^*\deleted{\text{\unsure{ens. de}}}$, où
\[
\underset{\substack{\Vert\\\operatorname{Aut}^-(V)\cup\{0\}}}{V'}\subset\operatorname{End}(V)
\]
est un sous-espace vectoriel de rg $2$ canonique, can. isom. à $V$, en associant à $\forall\ x\in V^*$ la transposition $\sigma_x\in\operatorname{Aut}^-(V)$.

2) $\operatorname{Aut}^+(V)=W^*$, où
\[
\underset{\substack{\Vert\\\operatorname{Aut}^+(V)\cup\{0\}}}{W}\subset\operatorname{End}(V)
\]
est un sous-espace vectoriel de rang $2$ canonique.

3) $V'$ et $W$ sont orthogonaux l'un de l'autre pour la forme trace, qui sur $V'$ se réduit à la forme bil. sym. unique qui est alternée et non dégénérée
\margin{$\begin{pmatrix}0&1\\1&0\end{pmatrix}$}
\[
(\operatorname{Tr}xy=\begin{cases}
0&\text{si }x=y,\text{ ou }x\text{ ou }y=0\\
1&\text{sinon, i.e. }x,y\in V'{}^*,\ x\ne y
\end{cases}),
\]
et sur $W$ est non alternée $(\operatorname{Tr}\rho^2=\operatorname{Tr}\rho^{-2}=1,\ \operatorname{Tr}1^2=0)$ et bien sûr non dégénérée.
\margin{$\begin{pmatrix}1&0\\0&1\end{pmatrix}$}

\margin{\unsure{cf.}}
4) $V'$ \deleted{\unsure{qui en}} tant que $\mathbb F_2(\Gamma)$-module est isom. à $V$ et irréductible, tandis que $W$ est une extension non scindée
\diagram{\begin{tikzcd}[ampersand replacement=\&]
0\arrow[r]\&\mathbb F_2\arrow[r]\&W\arrow[r,"\operatorname{Tr}"']\&\mathbb F_2\arrow[r]\&0\\
\&\lambda\arrow[r,mapsto]\&\lambda.1\&\&
\end{tikzcd}}
\margin{\uline{NB} $\mathbb F_2.1$ est son propre orthogonal}
\hfill\unsure{TSVP}
\scanpage{43}
% archive foliation: 42
5) On a
\[
\unsure{\mathfrak{sl}(V)|\mathfrak{gl}(V)}\overset{\mathrm{def}}{=}
\operatorname{Ker}(M\xrightarrow{\operatorname{Tr}}\mathbb F_2)
\unsure{=|\equiv}V'+\mathbb F_2.1_V
\]
\scanpage{44}
\mspage{II}
% archive foliation: 43
% additional faint pencil annotation, upper left: (25 v)
Structure octaédrale dans $\mathbb P^1$
\scanpage{45}
\mspage{1}
% archive foliation: 44
1. Opération de groupes $\simeq\mathbb D_2$ sur la droite projective, et configuration octaédrale.

Soit $V$ un groupe qui est un espace vectoriel de rang $2$ sur $\mathbb F_2$, et considérons \unsure{un|une} \deleted{droite pro} fibré en droites projectives $X$ sur un schéma de base $S$. On s'intéresse aux opérations, fidèles sur chaque fibre, de $V$ sur $X$, i.e. aux monomorphismes
\[
(1)\qquad V_S\lhook\joinrel\longrightarrow\underline{\operatorname{Aut}}_S(X)=G
\]
Choisissons une base $e_1,e_2$ de $V$ sur $\mathbb F_2$, alors la donnée d'un hom
\[
V_S\longrightarrow\underline{\operatorname{Aut}}_S(X)\qquad\text{i.e.}\quad V\longrightarrow\operatorname{Aut}_S(X)
\]
qui soit mono sur $\mathbb F_2e_1$ et sur $\mathbb F_2e_2$, équivaut à celle de deux \inserted{\unsure{autom.} d'ordre deux de $X$, ou encore de} ``axes'' ($=$ diviseurs étales de degré $2$) $\Delta_1,\Delta_2$ sur $X$, la condition de commutation étant que $\Delta_2$ est stable par $\sigma_1$, ou encore $\Delta_1$ \unsure{stble|stable} par $\sigma_2$.
\margin{on est en car. $\ne2$\quad $\Delta_1=X^{\sigma_1}$\quad $\Delta_2=X^{\sigma_2}$}
\deleted{Pour $\sigma_1$ donné, i.e. $\Delta_1$ donné, la donnée de $\sigma_2$ équivaut donc à la donnée d'une section de $X/\{1,\sigma_1\}\longrightarrow$}
\inserted{La condition que (1) soit un mono signifie \unsure{aussi} que l'on \unsure{a :}}
\[
(2)\qquad\Delta_1\cap\Delta_2=\varnothing\qquad\text{i.e.}\quad\Delta_2\subset X\unsure{-|\setminus}\Delta_1
\]
et pour $\sigma_1$ \inserted{i.e. $\Delta_1$} donné, la donnée de \deleted{$\unsure{\Delta_2}$} $\sigma_2$ i.e. de $\Delta_2$ équivaut à celle d'une section de \deleted{$\unsure{X\unsure{-|\setminus}\Delta_1}$ \unsure{sur}} $(X\unsure{-|\setminus}\Delta_1)/\{1,\sigma_1\}$ sur $S$ ($\Delta_2$ étant alors l'image inverse de cette section dans $X\unsure{-|\setminus}\Delta_1$). Considérons le troisième larron
\[
(3)\qquad e_3=e_1+e_2
\]
il lui correspond un axe $\Delta_3$, et $\Delta_1,\Delta_2,\Delta_3$ sont mutuellement disjoints.
\scanpage{46}
% archive foliation: 45
. La donnée de $\Delta_1$ est unique à conjugaison locale près - si on prend une trivialisation de \deleted{\unsure{donnée de $X$ triv.}} $\Delta_1$, \deleted{on peut prendre} la situation $(X,\Delta_1)$ i.e. \deleted{\unsure{on}} est isom. localement à $(\mathbb P^1_S,\{0,\infty\})$, $\sigma_1$ étant donné par $z\longmapsto-z$, donc $X\setminus\Delta_1$ étant \deleted{$\unsure{G}$} $\mathbb P^1_S\setminus\{0,\infty\}=G_{mS}$, et le quotient $(X\setminus\Delta_1)/\{1,\sigma_1\}$ s'identifiant à $G_{mS}$, l'hom de projection quotient étant
\diagram{\begin{tikzcd}[ampersand replacement=\&,row sep=large,column sep=large]
X\setminus\Delta_1\arrow[r]\arrow[d,phantom,"{\rotatebox{90}{$\simeq$}}" description]\&(X\setminus\Delta_1)/\{1,\sigma_1\}\arrow[d,phantom,"{\rotatebox{90}{$\simeq$}}" description]\\
G_{mS}\arrow[r,"\lambda\mapsto\lambda^2"]\&G_{mS}
\end{tikzcd}}
La donnée d'une section de $(X\setminus\Delta_1)/\{1,\sigma_1\}$ équivaut alors à celle d'une section $\lambda$ de $G_{mS}$ - et loc pour la top. étale, on peut poser $\lambda=\unsure{\zeta|\xi}^2$, \deleted{et quitte à modifier} \inserted{i.e. la section} \deleted{la} \inserted{$\lambda$} provient d'une section $\zeta$ de $X\setminus\Delta_1$ sur $S$. On peut (quitte à remplacer l'iso. $X\simeq\mathbb P^1_S$ par un autre, savoir en multipliant par $1/\zeta$) supposer que c'est la section $1$, de sorte que la situation \unsure{ét.-localt} est
\[
(4)\qquad\left\{\begin{array}{ll}
\Delta_1=\{0,\infty\},&\Delta_2=\{1,-1\}\\[1ex]
\sigma_1(z)=-z&\sigma_2(z)=\dfrac1z
\end{array}\right.
\]
(NB $\sigma_2$ est l'unique homographie qui \unsure{échange} $0$ et $\infty$, et fixe $1$ et $-1$). On trouve alors
\scanpage{47}
\mspage{2}
% archive foliation: 46
\[
(5)\qquad\sigma_3(z)=-\frac1z\qquad
\Delta_3=V(z^2+1)\quad(=\{i,-i\}\text{ si }\exists\text{ racine primitive de }\unsure{\epsilon-1|4\text{-}1}\text{ sur }S)
\]
On trouve donc

\uline{Proposition} Loc pour la top. étale, une représentation monomorphique de $\mathbb D_4=\mathbb F_2\times\mathbb F_2$ dans $\operatorname{Aut}(\mathbb P^1_S)$ (sur une base $S$ à car. résid. $\ne2$) est donnée par (4) (où $\sigma_1,\sigma_2$ sont les autos correspondant aux éléments de la base canonique $(e_1,e_2)$ de $\mathbb D_4$ sur $\mathbb F_2$) - qui donne bien une représentation monomorphique.

Posons
\[
(6)\qquad\Delta=\Delta_1\amalg\Delta_2\amalg\Delta_3
\qquad(=\{0,\infty,1,-1,i,-i\}\text{ dans le cas normalisé})
\]
on appelle $\Delta$ la \uline{configuration octaédrale} associée à l'opération de $V$ sur \deleted{$\unsure{\leftarrow}$} $X$ - par descente plate, elle est définie (comme diviseur étale de degré $6$ de $X/S$) dès qu'on a une opération monomorphique \inserted{(dans chaque fibre)} d'un groupe $V$ \inserted{$\simeq\unsure{\mathbb D_3|\mathbb D_2}$} sur $X$, ou plus généralement d'un schéma en groupes $\underline V$ sur $S$ loc. isom. au schéma en groupes constant. \deleted{\unsure{Alors}} De façon générale, on appelle \uline{configuration octaédrale} dans $X/S$ un diviseur étale de degré $6$ de $X/S$, qui soit
\scanpage{48}
% archive foliation: 47
loc. isomorphe à $\{0,\infty,1,-1,i,-i\}$ de $\mathbb P^1_S$, ou ce qui revient au $=$, qui soit loc. isom. à la configuration octaédrale associée à une op. \inserted{monom.} d'un $\underline V$ sur $X$. Nous allons voir maintenant que cette donnée est essentiellement unique.

Pour ceci, pour un diviseur octaédral \deleted{$\unsure{\Delta}$}
\[
\Delta\subset X
\]
considérons le schéma en groupes
\[
(7)\qquad\Gamma=\underline{\operatorname{Aut}}_S(X,\Delta),
\]
Il est immédiat (du fait que $\Delta$ est de degré $\geq3$) que $\Gamma$ \inserted{$\subset\underline{\mathfrak S}_{\Delta}=\underline{\operatorname{Aut}}_S(\Delta)$} est \uline{net} sur $S$, Je dis que $\Gamma$ est étale sur $S$, et loc. isom. au schéma en groupes $(\mathfrak S_4)_S$. Plus précisément, je vais définir sur $\Delta$ \inserted{(comme schéma de revêtements)} une structure d'\deleted{\unsure{oct.}} schéma en octaèdres combinatoires, de telle façon que l'on ait
\[
(8)\qquad\Gamma=\underline{\operatorname{Aut}}^+_{\mathrm{oct.}}(\Delta)
\qquad(\text{en tant que sous-schéma de }\underline{\operatorname{Aut}}_S(\Delta))
\]
Pour ceci, on \deleted{\unsure{se donne une}} \inserted{note qu'une} structure octaédrale sur un ens. à \unsure{six|dix} éléments est bien déterminée quand on connaît le groupe de ses automorphismes \inserted{directs} (en tant que sous-groupe de $\mathfrak S_\Delta$), car
\scanpage{49}
\mspage{3}
% archive foliation: 48
les ``diagonales'' $\Delta_1,\Delta_2,\Delta_3$ de $\Delta$ sont déterminées \inserted{(en tant que parties à deux éléments de $\Delta$)} \deleted{par la condition que leur stabilisateur soit} comme les classes pour la relation d'équivalence $R$ définie par l'application $s\longmapsto\Gamma_s$ ($s$ et $s'$ sont dans une $=$ classe ssi ils ont le $=$ groupe d'isotropie, qui est \inserted{cyclique} d'ordre $4$), et la structure polyédrale est déterminée en termes de cette décomposition de $\Delta$, les arêtes étant les \inserted{12} parties à deux éléments distinctes des $\Delta_i$, et les \inserted{8} faces étant les parties à trois éléments \uline{sections} de $\Delta\longrightarrow\Delta/R$.

Donc la structure octaédrale sur $\Delta$ revient à la structure définie par $R$, relation d'équivalence de degré $2$, ou encore par celle du groupe $\Gamma\ (\simeq\mathfrak S_4)$ opérant sur $\Delta$.

Cette procédure combinatoire montre que la structure octaédrale \unsure{annoncée} sur $\Delta\subset X$ est nécessairement unique (puisqu'on connaît son schéma en groupes d'autom.), ce qui montre que la question d'existence est locale (ét.), donc \unsure{OPS}
\[
X=\mathbb P^1_S,\qquad\Delta=\{0,\infty,1,-1,i,-i\}.
\]
Le \deleted{stabilisateur de $0,\infty$} \inserted{sous-groupe de} $\unsure{\Gamma}$ qui fixe $0$ et $\infty$ est isomorphe visiblement à
\[
\mu_4\deleted{\unsure{\simeq}}\simeq(\mathbb Z/4\mathbb Z)_S
\]
(qui contient $i$), \deleted{et par} et de $=$ (par raison d'homogénéité) pour $(1,-1)$, $(i,-i)$. Par contre, le
\scanpage{50}
% archive foliation: 49
groupe des autom. de $(X,\Delta)$ qui fixe $\infty$ et $-1$, donc \inserted{des autom.} de la forme $z\longmapsto az+b$ avec $-a+b=-1$ i.e.
\[
z\longmapsto\deleted{\unsure{a(z-1)+1}=az+(1-a)}\quad az+(a-1)
\]
est plus petit, en fait si l'autom. envisagé n'est pas l'identité, il doit induire une permutation sans pts fixes de $(0,1,i,-i)$, qui est donc d'ordre $2$ ou quatre, donc on a $a^4=1$ \deleted{ou $a^4=1$}.
\deleted{i.e. $a\in\unsure{\{1,i,-i\}|\{0,1,i,-i\}}$}
mais faisant $z=0$ on doit avoir
\[
a-1\in\{1,i,-i\}\qquad\text{i.e.}\quad a\in(2,1+i,1-i)
\]
faisant $z=1$ on trouve
\[
2a-1\in\{0,i,-i\}\qquad\text{i.e.}\quad a\in\left\{\frac12,\frac12(i+1),\frac12(1-i)\right\}
\]
\margin{donc $a^4=\pm4$ qui ne peut être $=1$ qu'en car. $3$ et $5$. \unsure{faux}, $a^4=-1$}
\margin{donc $a^4=1/4$ qui n'est pas ég. à $1$ (car. on est en car. $\ne2$) [\unsure{faux}.]}
faisant $z=i$, on trouve
\[
a\underbrace{(i+1)}_{1:\left(\frac{1-i}{2}\right)}-1\in\{0,1,-i\}
\qquad\text{i.e.}\quad
 a\in\left\{\deleted{\unsure{\frac1{i+1}}}\inserted{\unsure{\frac{1-i}{2}}},\deleted{\unsure{\frac2{i+1}}}\inserted{1-i},-i\right\}
\]
et de $=$, faisant $z=-i$, on trouve
\[
a\in\left\{\frac{1+i}{2},1+i,i\right\}.
\]
Donc l'\deleted{\unsure{a}}intersection des stabilisateurs de $\infty,-1$ est $\{1\}$, et par raison d'homogénéité il en est de $=$ pour toute partie à deux éléments de $\Delta$, distincte des $\Delta_1,\Delta_2,\Delta_3$. Donc on voit que la structure octaédrale $\Delta=\Delta_1\amalg\Delta_2\amalg\Delta_3$ est \unsure{reconnue|préservée} par $\Gamma$, et uniquement déterminée par le diviseur $\Delta$. Donc $\Gamma\subset\underline{\operatorname{Aut}}_{\mathrm{oct}}(\Delta)$, Pour montrer qu'on a isomorphisme, il suffit de noter que le groupe des autom. \inserted{directs} de la structure octaédrale est engendré par les trois sous-groupes cycliques d'ordre $4$ \deleted{\unsure{tous}} qu'il contient, correspondant à $\Delta_1,\Delta_2,\Delta_3$ -- or ceux-ci sont contenus dans $\Gamma$.
\margin{\unsure{et donc|et même} $\Gamma\subset\operatorname{Aut}_{\mathrm{oct}}^+(\Delta)$ \unsure{$=$}}
\scanpage{51}
\mspage{4}
% archive foliation: 50
Soit $\Delta$ un diviseur étale de degré $6$ dans $X/S$, muni d'une structure octaédrale i.e. d'une \deleted{\unsure{flèche}} \deleted{de degré $2$} \inserted{opposition involutive}
\[
\underline a_\Delta:\Delta\longrightarrow\Delta
\qquad\deleted{\Delta\longrightarrow\Delta'}
\]
\deleted{où $\Delta'$ est fini étale de degré $3$ sur $S$.}
Correspondant à une opération libre de $\{1,a_\Delta\}$ sur $\Delta/S$. On dit que cette structure octaédrale est \deleted{\uline{homogène}, \unsure{dans}} \inserted{\uline{régulière} dans} $X$, si le groupe
\[
\Gamma=\underline{\operatorname{Aut}}^+_{\mathrm{oct}}(\Delta)
\]
opère sur $X$, de façon compatible avec les op. sur $\Delta$ (cette op. de $\Gamma$ sur $X$ est alors unique).

\margin{Il faut un petit argument pour \unsure{montrer} que tout autom. de la structure oct. $\Delta$ induit par un autom. de $(X,\Delta)$ est direct : OPS que l'autom. fixe les points de $\{0,\infty\}=\Delta_1$, et on se trouve dans les autom. de \unsure{``ceinture''} par $\mu_4$, qui sont bien directs\ldots}

\uline{Pour} $\Delta$ \inserted{diviseur} fixé, on voit qu'il y a au plus une structure octaédrale sur $\Delta$ qui soit \uline{régulière}, et il y en a une ssi $\Delta$ est une ``configuration octaédrale''. Si $\Delta$ est une telle configuration, on vient de voir \inserted{déjà} qu'elle est munie d'une \deleted{unique} structure octaédrale régulière (et \uline{tout} autom. de $X$ qui \deleted{fixe} \inserted{conserve} $\Delta$ \deleted{fixe} \inserted{conserve} cette structure octaédrale). Inversement, si $\Delta$ est muni d'une str. oct. rég, alors le sous-groupe $\underline V$ de $\underline\Gamma=\underline{\operatorname{Aut}}_{\mathrm{oct}}\deleted{\unsure{\mathrm{oc}}}(\Delta)$ inv. abélien non trivial est un groupe loc. isom. à $\mathbb D_{2S}$, qui opère fidèlement sur $X$, et on voit aussitôt que le

\margin{Si on a une configuration octaédrale régulière \deleted{\unsure{$\Delta$}} dans un \deleted{\unsure{fibré}} \inserted{\unsure{en} droites projectives} $X$, alors le revêtement des orientations de \deleted{\unsure{$\Delta$}} est can. isom. au revêtement $\mu_4^*\unsure{=}\operatorname{Spec}\mathcal O_S[i]/(i^2+1)$ de $S$\ldots}
\scanpage{52}
% archive foliation: 51
configuration \deleted{\unsure{icosaèdre}} \inserted{octaédrale} associée est bien celle de départ\ldots

\uline{Théorème}\quad Soit $X$ un fibré en droites projectives sur un schéma $S$, de car. résiduelles $\ne2$. Considérons le \deleted{\unsure{gr.}} schéma en groupes $G=\underline{\operatorname{Aut}}_S(X)$ (\uline{NB} la donnée de $X/S$ équivaut à celle de la ``forme'' $G$ de $(\operatorname{PGL})_S$ sur $S$), et les ensembles suivants :

 a) L'ens $A(X)$ des sous-schémas en groupes $V$ de $G$, loc isom. à $(\mathbb D_4)_S$.

 b) L'ens \deleted{\unsure{de}} $B(X)$ des sous-schémas en groupes \deleted{\unsure{$G$}} $\Gamma$ de $G$, loc. isom. à $(\mathfrak S_4)_S$.

 c) L'ens des configurations icosaédrales régulières dans $X$, i.e. des diviseurs étales $\Delta$ de $X/S$, de degré $6$, munis d'une opposition $\underline a_\Delta:\Delta-\Delta$ qui \deleted{\unsure{définit} des pts} fixes sur chaque fibre géom., et tel que $\underline{\operatorname{Aut}}^+_{S\mathrm{oct}}(\Delta)$ opère sur $X$ de façon compatible avec son opération sur $\Delta$.

On a alors un système transitif de bijections canoniques entre ces trois ensembles, \unsure{défini} ainsi !
\scanpage{53}
\mspage{5}
% archive foliation: 52
\[
 a)\longrightarrow b)\qquad \Gamma=\underline{\operatorname{Norm}}_G\underline V
\]
$b)\longrightarrow a)$\quad $\underline V$ est l'unique sous-schéma en groupes \inserted{distingué} fini étale de $\Gamma$, \deleted{\unsure{q}} de degré $4$, \deleted{distingué}.
\margin{\uline{NB}\quad $\Gamma/V\simeq\underline{\operatorname{Aut}}_{S\mathrm{gr}}(V)\simeq\underline{\operatorname{Aut}}_S(V^*)$}
\[
 c)\longrightarrow b)\qquad
 \Gamma=\underline{\operatorname{Aut}}_S(X,\Delta)\xleftarrow{\sim}\underline{\operatorname{Aut}}^+_{S\text{-}\mathrm{oct}}(\Delta)
\]
$\unsure{a})\longrightarrow c)$\quad \deleted{$\underline V$ est} Toute section $\inserted{\sigma\in}V^*=V\setminus(\text{section unité})$ définit une action \inserted{fidèle} de \deleted{\unsure{ce}} $(\mathbb Z/2\mathbb Z)_S$ sur $X$, donc un ``axe'' $\Delta_\sigma$ de $X$, et $\Delta$ est \deleted{la réuni\unsure{on}} (loc. ét.) le \deleted{\unsure{sché}} diviseur somme de ses axes\ldots

\uline{Corollaire 1} Pour un diviseur étale de degré $6$ $\Delta$ dans $X$, il existe au plus une structure icosaédrale régulière ; \deleted{dont le groupe des autom.} \inserted{S'il en existe une, le} \inserted{schéma en} groupe de ses automorphismes \inserted{directs} est isom.
\[
\underline{\operatorname{Aut}}_S^{\unsure{+}}(X,\Delta).
\]

\uline{Corollaire 2}\quad La catégorie des couples $(X,\Delta)$ \inserted{munis d'un iso $\underline{\operatorname{Or}}\Delta/S\simeq\mu_{4S}^*$ (on dit qu'ils sont $\mu_4^*$-orientés)} sur $S$, (i.e. $(X,\underline V)$, \deleted{\unsure{orientés}} $(X,\underline\Gamma)$, \inserted{$\simeq$ \deleted{\unsure{$V$}} octaèdres combinatoires $\Delta\ldots$}) est équivalente à celle des \deleted{\unsure{sur}} $S$, ou encore : celle des revêtements finis étales \deleted{$\unsure{\Delta}$} \inserted{$\widetilde F$} de degré $4$ de $S$, en associant à \deleted{$\unsure{H}$} $(X,\Delta)$ le revêtt des quatre ``faces doubles antipodiques'' de la structure octaédrale $\Delta$, ou ce qui revient au $=$, celui des sous-$3$-groupes de Sylow de \deleted{$\unsure{\mathfrak S_4}$}
\[
\underline{\operatorname{Aut}}^+_{S\mathrm{oct}}(\Delta)\simeq\underline{\operatorname{Aut}}_S(X,\Delta).
\]
La structure-type${}^*$ correspond \deleted{\unsure{au choix de}} $X\subset\mathbb P^2_{\mathbb Z[\frac12]}$, $X=V(x^2+y^2+z^2=0)$, avec l'op. de \deleted{\unsure{$G$}} $\{\pm\}^3/\mathrm{diag}\simeq V_{\{0,1,\infty\}}(\mathbb F_2)$ donnée par les chgts de signes en $x,y,z$${}^*$.

\margin{${}^*$ La \unsure{biface|face} distinguée est formée des deux faces $\{(0,1,i),(i,0,1),(1,i,0)\}$ et $\{(0,1,-i),(-i,0,1),(1,-i,0)\}$.}
\margin{\uline{NB} Ici $X$ est totalement anisotrope sur $\mathbb R$ i.e. n'a pas de pts réels sur $\mathbb R$ -- le revêtement galoisien associé étant de groupe $\unsure{V}_{\{0,1,\infty\}}(\mathbb F_2)$ de $\mathbb P^1_{\unsure{\mathbb C}}$, ramifié \unsure{seulement} en $0,1,\infty$\ldots}
\scanpage{54}
% archive foliation: 53
2. \uline{Courbes relatives de type} $(0,4)$.

Soit $X$ un fibré en droites projectives sur $S$, et
\[
(9)\qquad\underline I\subset X
\]
un diviseur relatif, étale de degré $4$. Considérons
\[
(10)\qquad\underline{\mathfrak S}_{\underline I}=\underline{\operatorname{Aut}}_S(I)
\]
qui est un schéma en groupes fini étale sur $S$, loc isom. à $(\mathfrak S_4)_S$. Il lui correspond un unique sous-groupe \inserted{distingué} $\underline V_{\underline I}$, loc. isom. à $(\mathbb D_2)_S$, \deleted{de telle faç\unsure{on}} Mais $\underline I$ est un torseur sous $\underline V_I$.

\uline{Proposition}\quad L'opération de $\underline V_{\underline I}$ sur $\underline I$ se prolonge (de façon évidemment unique) en une opération de $\underline V_{\underline I}$ sur $X$.

\uline{NB} La donnée de $\underline I$, fini étale de degré $4$ sur $S$, équivaut à celle de $\underline V=\underline V_{\underline I}$ sur $S$ (ou encore de $\underline V_{\underline I}^*=\underline J_{\underline I}$, fini étale de degré $3$ sur $S$), plus la donnée d'un torseur $\underline I$ sous $\underline V^*$. Plonger alors $\underline I$ dans un fibré en droites projectives $X$, revient à faire opérer fidèlement $\underline V$ sur $X$ (chose étudiée en § 1), et de plus, de \unsure{se}
\scanpage{55}
\mspage{6}
% archive foliation: 54
donner une section de $(X\setminus\Delta)/\underline V$ \deleted{\unsure{de}}, en prenant \inserted{pour $I$} l'image inverse de cette section dans $X\setminus\Delta$.

Mais la donnée d'une structure octaédrale équivaut à celle d'un \deleted{\unsure{$I$}} rev. étale de degré quatre de $S$ (à ne pas confondre avec le $\underline I$ de tantôt, appelons le $\underline K$), ou encore, à celle de

 a) un schéma en groupes $V$ sur $S$, loc. isom. à $(\mathbb D_2)_S$ -- ou encore un revêtement fini étale $\underline J\ (\unsure{\Leftarrow|=}V^*)$ de degré $3$ de $S$, ou \deleted{\unsure{de}} enfin, à celle d'un \deleted{\unsure{ce}} fibré en droites projectives $\bar\Sigma=\Sigma_{\underline J}$ sur $S$, muni d'un diviseur \inserted{relatif} étale $D_J$ de degré $3$ (\unsure{s. isom.} \deleted{\unsure{lui}} à $\underline J$)

 b) d'un torseur $\underline K$ sous $\underline V$, ou ce qui revient au $=$, d'un revêtement \deleted{galoisien} \inserted{\uline{principal}} $X^*$ de $\Sigma^*_{\underline J}=\Sigma_{\underline J}\setminus\deleted{\unsure{J}}D_{\underline J}$ de groupe \deleted{$\unsure{H^0(S,\mathscr H^1(\ ,V))}$} \inserted{$(\underline V_{\underline J})_{\Sigma^*_{\underline J}}$}, \deleted{géométriquement connexe, de degré \unsure{fibre}} \inserted{correspondant à $\unsure{\xi}$ un élément de
\[
H^0(S,\mathscr H^1(\Sigma^*_{\underline J}/S,\underline V))\simeq\operatorname{Hom}(\underline V_{\underline J},\underline V)\deleted{\unsure{\simeq\mathfrak S_3}}
\]
qui est l'identité $\underline V_{\underline J}\longrightarrow\underline V=\underline V_{\underline J}$} -- la correspondance entre $\underline K$ et $X$ étant donnée par la formule
\[
(11)\qquad X^*=\deleted{\unsure{X_{\underline J,0}\times}}\quad \underline K\wedge_V\unsure{\widetilde\Sigma}_{\underline J,2}^*,
\]
où $\unsure{\widetilde\Sigma}_{\underline J,2}^*$ est le revêtement principal dé

\margin{\uline{NB} Le grpt. \unsure{local des} $(\pi_1)_{\mathrm{ab}/2}$ de $\Sigma^*_{\underline J}$ sur $S$ est can. isom. à $\underline V_{\underline J}$\ldots}
\scanpage{56}
% archive foliation: 55
$\Sigma^*_{\underline J}$ défini \deleted{\unsure{par}} à partir de la \deleted{fibré} conique \deleted{\unsure{def}} d'équation homogène
\[
(12)\qquad\sum_{i\in\underline J}y_i^2=0
\]
sur laquelle le groupe $V_{\underline J}$ \inserted{$\simeq\mathbb F_2^{\underline J}/\mathbb F_2$} opère \deleted{\unsure{par}} par changements de signes sur les $x_i$, \unsure{a} comme quotient $\Sigma_{\underline J}$, défini par l'équation homogène
\[
(13)\qquad\sum x_i=0.
\]
\margin{description valable loc. (ét.) ; sinon, il faut prendre dans $\mathbb P(p_*(\mathcal O_{\underline J}))$, \deleted{\unsure{de}} la conique d'équation $Q=0$, où $Q(u\deleted{\unsure{u}})=\operatorname{Tr}u^2$ (NB $p:\underline J\longrightarrow S$ le morphisme structural).}
\margin{description valable loc. ét. -- la description est par l'hyperplan (13) de $\mathbb P(p_*(\mathcal O_{\underline J}))$, d'équation $\unsure{L(u)}=0$, où $L(u)=\operatorname{Tr}u$.}

Donc La donnée de $(X,\underline I)$ équivaut à celle de \deleted{\unsure{$\Sigma$}} $\underline J$ i.e. de $\Sigma_{\underline J}$, d'un torseur $\underline K$ sous $V_{\underline J}$ i.e. d'un revêtement $X$ de $\Sigma_{\underline J}$ \deleted{ramifié} de groupe $V_{\underline J}$ (admissible au sens qu'il définit l'\uline{identité} $V_{\underline J}\longrightarrow V_{\underline J}$), \uline{et} d'une section $\unsure{t|x}$ de $\Sigma^*_{\underline J}$.

Quand $\underline I$ est \inserted{fixé et} de la forme $I_S$, avec $I$ ens. à quatre éléments fixe (à fins d'épinglage), on a $\underline J\simeq J_S$, et $\Sigma^*_{\underline J}=\bar\Sigma_J$ est aussi fixé. La donnée de $(X,I_S\hookrightarrow X)$ équivaut alors à celle de la section $x$ de $\Sigma^*_J$, servant de ``pt référence'', le revêtement $X$ de $\bar\Sigma^*_J$ de groupe $(V_{\underline J})_S=(V)_{(J_S)}=\underline V$ étant
\margin{[L'hypothèse que $\underline I$ soit constant est en fait superflue]}

\margin{Relations entre $\underline I,\underline K,x$ :
\[
(14)\quad\underline I\simeq\underline K\wedge_{\underline V}\bigl(\unsure{\widetilde\Sigma}_{\underline J,2}^*|x\bigr).
\]
Ainsi, pour $\underline J$ et $x\in\Gamma(\Sigma^*_{\underline J}/S)$ fixés, $\underline I$ et $\underline K$ (ou encore, $\underline I$ et $X$) se déterminent mutuellement\ldots On a trois revêtements étales de degré $4$ de $S$, $\underline I$, $\underline J\amalg S$, et $\underline K$\ldots}
\scanpage{57}
\mspage{7}
% archive foliation: 56
\unsure{déterminé} par la \deleted{\unsure{f}} structure supplémentaire
\[
(15)\qquad X|t\xrightarrow{\sim}I_S\qquad(V_S\text{-isomorphisme})
\]
qui, en termes du $V_S$-torseur $\underline K$ qui est censé définir $X$, s'écrit
\[
\deleted{(16)\qquad\underline K\simeq(I_S)\wedge_{\unsure{V_S}}\bigl(\unsure{\widetilde\Sigma}_{J,2}^*|t\bigr)}
\]
\[
(16)\qquad\underline K\simeq(I_S)\wedge_{\underline V}(\Sigma_{J,2}^*|t)
\qquad\underline V=(V_J)_S.
\]
\scanpage{58}
% archive foliation: 57
3: Relations avec courbes elliptiques:
\inserted{\uline{le morphisme} $f_2^t:\Sigma_{\underline J}\longrightarrow\bar\Sigma_{\underline J}$.}

On suppose tjrs les car. rés. de $S\ne2$. Soient $E_0$ une courbe elliptique sur $S$, $E$ un torseur sous $E_0$ (de sorte que $E_0\simeq\underline{\operatorname{Pic}}^0_{E/S}$).

On suppose donné dans $E$ un $S$-autom
\[
\tag{17}\sigma:E\longrightarrow E.
\]
tel que
\[
\tag{18}\sigma_{E_0}^{\unsure{u|h}}=-\operatorname{id}_{E_0}
\]
ce qui implique que loc. (ét.), identifiant $E$ à $E_0$ via une section origine de $E$, on aura
\[
\sigma x=a-x\qquad a\text{ section convenable}
\]
ce qui est d'ailleurs intrinsèque si on interprète $a$ comme une section de $E\wedge_{E_0}E\simeq\underline{\operatorname{Pic}}^2_{E/S}$.

On aura donc
\[
\sigma^2=\operatorname{id}_E.
\]
Soit \deleted{$\unsure{\phi|I}$}
\[
\tag{19}\left\{\begin{array}{ll}
\underline I=E^\sigma&\text{sous-schéma des pts fixes de }\sigma\\
\underline I\subset E&\\[-4pt]
\Vert&\\[-4pt]
\{x\in E\mid 2x=a\}&
\end{array}\right.
\]
on voit que $\underline I$ est un torseur sous le schéma en groupes
\[
\tag{20}\underline V={}_2E_0=\operatorname{Ker}(2\operatorname{id}_{E_0}:E_0-E_0),
\]
qui est justement un groupe \inserted{sur $S$} loc. isom. à $\mathbb D_{2S}$.

La donnée de $(E,\sigma)$, où $E$ est un schéma
\scanpage{59}\mspage{8}
% archive foliation: 58
propre et lisse sur $S$ à fibres géom. connexes de dim. $1$, genre $1$, et $\sigma$ un autom induisant $-\operatorname{id}$ sur $\underline{\operatorname{Pic}}^0_{E/S}$, \uline{équivaut} à celle \inserted{de $(E,a)$, où $a$ section de $\underline{\operatorname{Pic}}^2_{E/S}$ $(\simeq E\wedge_{E_0}E)$, ou encore à celle} \deleted{de} $(E_0,\underline I)$, où $E_0$ est une courbe elliptique relative sur $S$, et où $\underline I$ est un torseur sous \deleted{${}_2E_0\overset{\mathrm{déf}}=$} $\underline V\overset{\mathrm{déf}}={}_2(E_0)$. On aura
\[
\tag{21}\left\{\begin{array}{l}
E=\underline I\wedge_{\underline V}E_0\\
\sigma\text{ déduit de la symétrie }x\mapsto-x\text{ de }E_0\\
\qquad\text{via la formule précédente}
\end{array}\right.
\]
On aura une opération de $\underline V\unsure{\Leftarrow|=}{}_2E_0\subset E_0$ sur $E$, et un isom can.
\[
\tag{22}\begin{array}{ccc}
E/\underline V&\simeq&E_0\\
\rotatebox{90}{$\simeq$}&&\\
E\wedge_{E_0}E\simeq\underline{\operatorname{Pic}}^2_{E/S}&&
\end{array}
\]
\margin{\uline{NB} L'opération de $\underline V$ sur $E$ commute à l'action de $\sigma$}
isom qui commute aux symétries $\sigma$ sur $E$, $\sigma_0:x\mapsto-x$ sur $E_0$.

Considérons les schémas quotients
\[
\tag{23}X=E/\{1,\sigma\},\qquad\Sigma=E_0/\{1,\sigma_0\}
\]
ce sont des fibrés en \deleted{\unsure{droites}} \inserted{droites} projectives,
\deleted{\unsure{et on a les}} $V$ opère \inserted{(monomorphiquement)} sur $X$ \deleted{et sur $\Sigma$} via son \deleted{\unsure{(24)}} action sur $E$, et on aura
\[
\tag{24}X/V\simeq\Sigma,
\]
d'ailleurs le diviseur de ramification pour
\[
\tag{25}E\longrightarrow X\qquad(\text{rev. quadratique})
\]
\scanpage{60}
% archive foliation: 59
est $I$, et on a un \deleted{\unsure{mono}}\uline{morphisme}
\[
\tag{26}I\hookrightarrow X
\]
\deleted{stable avec l'\unsure{action}}
compatible avec l'action de $\underline V$, d'où on déduit facilement que c'est un \uline{monomorphisme} (car on vérifie que $I$ s'envoie dans l'ouvert $\underset{(27)}{X^*}$ de $X$ où l'action de $\underline V$ est libre, cf plus bas).

Ainsi on est dans la situation du § 2, qui implique que $\Sigma\simeq\bar\Sigma_{\underline J}$, où $\underline J=\underline V^*$, que $I$ est l'image inverse d'une section bien déterminée $t$ de $\Sigma_{\underline J}^*=\Sigma_{\underline J}\setminus D_{\underline J}$, l'image inverse $\Delta$ de $\underline J$ dans $X$ étant \deleted{\unsure{une}} la configuration icosaédrale dans $X$, canoniquement associée à l'action donnée de $\underline V$ sur \deleted{le} fibré en droites projectives $\underline V$ sur $X$.

En termes de l'action de $\underline{\unsure{V|V^*}}$ sur $X$, $\underline\Delta$ se récupère comme le diviseur de ramification pour cette action -- ensemblistement, il correspond donc aux pts (à val. dans des corps alg. clos) dont le stabilisateur dans $V$ est $\ne0$, i.e. il correspond à des pts $x\in E$ tels qu'il existe $\zeta\in V={}_2E_0$, $\zeta\ne0$, avec
\[
 x+\zeta\in\{x,\sigma x=a-x\}\quad\text{i.e.}
\]
$x+\zeta=a-x$ i.e. $2x+\zeta=a$ d'où $4x=2a$ \deleted{i.e. $\unsure{x\in I\wedge_4E_0|x\in I\wedge_2E_0}$} \inserted{mais $2x\ne a$}
\scanpage{61}\mspage{9}
% archive foliation: 60
$x$ pt de $I\wedge_{{}_2E_0}({}_4E_0)\setminus I$

On trouve ainsi
\[
\tag{27}\Delta=
\underbrace{\left(\underbrace{\underbrace{I\wedge_{\underline M_0}\underline M}_{\text{degré }16}\setminus\underbrace{I}_{\text{degré }4}}_{\text{degré }12}\right)/\{1,\sigma\}}_{\text{degré }2}
\]
où on a posé maintenant
\[
\tag{28}\left\{\begin{array}{l}
\underline M\simeq{}_4E_0\quad\text{schéma en groupes loc. isom. à }(\mathbb Z/4\mathbb Z)^2_S\\
\underline M_0\simeq\underline M\otimes_{\mathbb Z/4\mathbb Z}\mathbb Z/2\simeq2\underline M\simeq{}_2E_0=\underline V
\end{array}\right.
\]
Ainsi \inserted{(dans le cas $E=E_0$, disons)} les pts d'ordre $4$ de $E_0$ se groupent par $6$ paires d'éléments symétriques, qui \unsure{fournit|fournissent} la configuration octaédrale $\Delta$ dans $X$.

On trouve ainsi le diagramme (presque cartésien)
\[
\tag{29}
\diagram{\begin{tikzcd}[ampersand replacement=\&,column sep=large,row sep=huge]
\underline I\arrow[r,phantom,"\subset"]\&E\arrow[r,"\text{étale}"]\arrow[d,"\substack{\text{ramifié en}\\\text{\unsure{dessus} des pts}\\\text{de }\underline I}"]\&E_0\;(\simeq E/\underline M_0)\arrow[d,"\substack{\text{ramifié en dessus}\\\text{des pts de }D_J\cup t(S)}"]\\
\underline I\arrow[r,hook]\&X\arrow[r,"\substack{\text{ramifié en}\\\text{dessus des}\\\text{pts de }D_J}"]\&\Sigma_J\\
\&\Delta\arrow[u,phantom,"\rotatebox{90}{$\subset$}"]\arrow[r,dashed]\&D_J\arrow[u,phantom,"\rotatebox{90}{$\subset$}"]
\end{tikzcd}}
\]
qui est cartésien au dessus de $\bar\Sigma_J^*=\bar\Sigma_J\setminus D_J$
\scanpage{62}
% archive foliation: 61
(au dessus duquel $X^*$ est étale \deleted{\unsure{galois}} principal de groupe $\underline V$), tout comme $E$ sur $E_0$) mais non au voisinage des pts de $D_J\subset\Sigma_J$ (vu que $E_0/\Sigma_J$ y est ramifié, mais non $E/X$ en les points de $\Delta=X\mid D_J$, or si c'était cartésien il en résulterait que $E_0/\Sigma_J$ est étale sur les pts de $D_J$ aussi, par descente plate\ldots). \unsure{Quand|Quant} $=$ il doit être facile de vérifier que le diagramme total se reconstitue à partir de $E_0$ ($\mapsto$ \inserted{$\underline J$ et} $\Sigma=\Sigma_{\underline J}$, et $E_0\longrightarrow\Sigma$, \inserted{et de section $t$ de $\Sigma^*$}) et du revêtement \deleted{principal} $X$ de $\Sigma$ avec groupe \inserted{\unsure{dég.}} $\underline V$: la donnée de $(E,\sigma)$ i.e. de $(E_0,I)$ \deleted{doit être} \inserted{est} équivalente aussi à $E_0$ et du revêtement $X$ (d'où en fait \deleted{$\underline I$}\inserted{$\underline K$} et $\underline I\simeq\underline K\wedge(X\mid t)\unsure{'|^{\circ}}\ldots$)

Cela montre que finalement $(E,\sigma)$ i.e. le $I$ est un peu de la poudre aux yeux, on va prendre pour le reste $E=E_0$, $\sigma=\sigma_0$, donc $I=\underline V={}_2E_0=M_0$, torseur trivial. On a donc \deleted{$\unsure{\Sigma|I}$} $X=\Sigma_J$ et on trouve donc
\scanpage{63}\mspage{10}
% archive foliation: 62
une application canonique
\[
\tag{30}\Sigma_{\underline J}\xrightarrow{f_{\unsure{2}}^t}\Sigma_{\underline J}
\]
\margin{c'est l'application déduite par passage aux quotients de $2\operatorname{id}_{E_0}:E_0-E_0$}
qu'on doit pouvoir se définir en termes de $\underline J$, et de la seule section $t$ de $\Sigma_{\underline J}$ (ce qui en fait détermine $E_0$, à isom \inserted{(non unique)} près loc. pour la top. étale\ldots). On a i.e.
\margin{il faudrait pouvoir \unsure{donner une} description intrinsèque de $\unsure{f_2^t|f_i^t}$, sans faire intervenir des courbes elliptiques}
\deleted{$\unsure{I}=\underline J\amalg$}
\[
\tag{31}\underline I_t=\underline J\amalg t(S)
\]
\margin{\uline{NB} comme $S$-schéma, $\underline I_t$ ne dépend pas de $t$, mod isom. can.}
et \deleted{$\unsure{E_0}$} $\underline V$ peut se décrire comme $V_{\underline I}\simeq V_{\underline J}$, opérant sur $\underline I$ de la façon canonique, cette opération se prolongeant en une opération sur $\Sigma_{\underline J}$ (dépendant de la section $t$ de $\Sigma_{\underline J}^*$). Le \deleted{\unsure{applic}} morphisme $f_t$ est défini sans ambiguïté par les conditions suivantes (où on suppose pour simplifier que $\underline J=J_S$, $J$ \unsure{ens|un} \unsure{à|:} trois éléments):
\[
\tag{32}\left\{\begin{array}{l}
\forall i\in J,\ \text{si }\sigma_i^t\in\operatorname{Aut}_S(\Sigma_J)\text{ est l'autom. qui échange}\\
t\text{ et }i,\text{ et }j\text{ et }k\quad(\text{où }\{i,j,k\}=J),\text{ et}\\
\deleted{\text{alors }\unsure{\text{si }(\Sigma_J)^{\sigma_i}}}\quad\Delta_i^t=(\Sigma_J)^{\sigma_i^t},\text{ on a}\\
f(\Delta_i^t)=\text{section }\xi_i\text{ de }D_{\underline J}\simeq J_S
\end{array}\right.
\]
et cela doit impliquer automatiquement
\[
\tag{33}f_t(\unsure{I_S|I_t}=J_S\cup t(S))\subset\deleted{\text{section}}t(S)
\]
\scanpage{64}
% archive foliation: 63
Faisons les calculs pour $J=\{0,1,\infty\}\subset\mathbb P^1_S$, on trouve facilement:
\[
\tag{34}\left\{\begin{array}{ll}
\sigma_0^t(z)=\dfrac{z-t}{z-1}&(\text{éch. }0\text{ et }t,\ 1\text{ et }\infty)\\
\sigma_1^t(z)=t/z&(\text{éch. }1\text{ et }t,\ 0\text{ et }\infty)\\
\sigma_\infty^t(z)=t\dfrac{z-1}{z-t}&(\text{éch. }\infty\text{ et }t,\ 0\text{ et }1)
\end{array}\right.
\]
\margin{NB $\sigma_\infty^t=\sigma_0^t\sigma_1^t$}
On trouve donc
\[
\tag{35}\left\{\begin{array}{lll}
\Delta_0^t:&\text{solutions de }z^2-2z+t=0,&\Delta_0^t=\{1\pm\sqrt{1-t}\}\\
\Delta_1^t:&\text{\rule{3cm}{0.4pt}}\quad z^2-t=0&\Delta_1^t=\{\pm\sqrt t\}\\
\Delta_\infty^t:&\text{\rule{3cm}{0.4pt}}\quad z^2-2tz+t=0&\Delta_\infty^t=\{t\pm\sqrt{t(1-t)}\}
\end{array}\right.
\]
et on doit avoir
\[
\tag{36}f_2^t:\quad\begin{array}{l}
\Delta_0^t\longrightarrow0\\
\Delta_1^t\longrightarrow1\\
\Delta_\infty^t\longrightarrow\infty
\end{array}\qquad\uline{\text{avec multiplicité }2}
\]
ce qui nous donne (puisque les zéros et pôles de $f_2^t$ de degré quatre sont connus)
\[
f_2^t(z)=c_t\left(\frac{z^2-2z+t}{z^2-2tz+t}\right)^2\qquad c_t\text{ section de }\underline O_S^*,
\]
dont la valeur s'obtient en écrivant que $\Delta_1^t\longrightarrow1$. Mais si $\zeta$ est une section de $\Delta_1^t$ i.e. $\zeta^2=t$, on a
\margin{$\begin{aligned}
\zeta^2\deleted{{}+t}-2\zeta+t&=2(t-\zeta)\\
\zeta^2-2t\zeta+t&=2t(1-\zeta)
\end{aligned}$}
\[
f_2^t(\zeta)=c_t\left(\frac{2(t-\zeta)}{2t(1-\zeta)}\right)^2
=c_t/t^2\ \deleted{\frac{\zeta^2-2t\zeta+t^2}{\zeta^2-2\zeta+1}}
=\left(\frac{t-\zeta}{1-\zeta}\right)^2
\]
Or
\[
\deleted{{}=\frac{c_t}{t^2}\frac{t^2+t-2t\zeta}{1+t-2\zeta}}\quad-
\quad\frac{t-\zeta}{1-\zeta}
=\frac{(t-\zeta)(1+\zeta)}{(1-\zeta)(1+\zeta)}
=\frac{-\zeta^2+(t-1)\zeta+t}{1-\zeta^2}
=\frac{(t-1)\zeta}{1-t}=-\zeta
\]
donc
\[
f_2^t(\zeta)=\frac{c_t}{t^2}\zeta^2=\frac{c_t t}{t^2}=\frac{c_t}{t}\qquad\text{d'où}\quad c_t=t
\]
\scanpage{65}\mspage{11}
% archive foliation: 64
et par suite
\[
\tag{37}f_2^t(z)=t\left(\frac{z^2-2z+t}{z^2-2tz+t}\right)^2
\]
Il reste à vérifier que les valeurs de cette fonction pour $z=0,1,\infty$ ou $t$, sont égales à $t$, i.e. que celle de $\dfrac{z^2-2z+t}{z^2-2tz+t}$ est égale à $\pm1$. Pour $0,\infty,1$ on trouve $1$, pour $z=t$ on trouve $\dfrac{t^2-t}{t-t^2}=-1$, OK.

\uline{NB} En prenant le \unsure{n.ième} itéré de $f_2^t$, on trouve le morphisme $f_{2^n}^t$, qui correspond par passage au quotient à $2^n\operatorname{id}_{E_0}$. L'image inverse de la section $t$ par $f_{2^n}^t=(f_2^t)^{\circ n}$ est l'image par $p:E_0\longrightarrow\Sigma_J\simeq E_0/\{1,\sigma_{E_0}\}$ du \deleted{\unsure{l'union}} \inserted{sous}-schéma \inserted{en groupes} ${}_{2^n}E_0$ des points annulés par $2^n$.

:
\scanpage{66}
% archive foliation: 65
4. Détermination de $M_{11}(4)$.

Soit $E$ une courbe elliptique sur $S$ (à car. résiduelles $\ne2$), considérons
\[
\tag{38}M={}_4E\qquad\begin{array}{l}
\text{schéma en groupes sur }S\\
\text{loc. isom. à }(\mathbb Z/4\mathbb Z)^2,\\
\text{i.e. }\deleted{\text{syst. local}}\text{ faisceau abélien}\\
\text{sur }S_{\text{ét}}\text{ qui est un }(\mathbb Z/4\mathbb Z)_S\\
\text{module loc. libre de rg }2.
\end{array}
\]
de sorte qu'on a
\deleted{(39)}
\[
\tag{\unsure{39|34}}\diagram{\begin{tikzcd}[ampersand replacement=\&]
0\arrow[r]\&M_0\arrow[r]\&M\arrow[r,"\mathrm{can}"]\&M_0\arrow[r]\&0\\
\&M\arrow[u,"\mathrm{can}"]\arrow[ur,"2\operatorname{id}_M"']\&\&\&
\end{tikzcd}}
\]
où
\[
M_0=M\otimes_{\mathbb Z/4\mathbb Z}\mathbb F_2\simeq{}_2E.
\]
La donnée d'un $M$ équivaut à celle d'un $M_0$ (i.e. d'un $\underline J$ \deleted{\unsure{rev.}} étale de degré $3$ de $S$ \deleted{$\unsure{\underline J}$}
\[
\tag{40}\underline J\simeq M_0^*,\qquad M_0=V_{\underline J}(\mathbb F_2)\simeq\mathbb F_2^{\underline J}/\underset{\mathrm{diag}}{\underline{\mathbb F_2}}),
\]
\uline{et} d'un torseur \inserted{$P$} sous
\[
\tag{41}\unsure{N|V}_{\underline J}(M_0)\simeq\underline{\operatorname{End}}_{\mathbb F_2,S}(M_0)\simeq M_0\oplus\mathbb F_2^{\underline\omega}
\]
\margin{$\underline\omega$ est le rev. étale de degré $2$ ``des orientations'' associé au rev. $\underline J$ de degré $3$ de $S$}
Un ``épinglage'' de $M={}_4E$ (i.e. \deleted{d'une} la donnée d'un iso
\[
\tag{42}(\mathbb Z/4\mathbb Z)^2_S\simeq M\;(={}_4E)
\]
équivaut à celle d'une \inserted{trivialisation de $\underline J$, plus d'une} section \unsure{(trivialisé)}-$\unsure{\underline P|\underline E}$
\scanpage{67}\mspage{(12)}
% archive foliation: 66
\deleted{plus celle d'une}
de ce torseur, \deleted{qui} \inserted{Ce torseur} peut aussi s'interpréter ainsi: la structure d'extension (39) \deleted{comme le torseur des sections} définit la structure d'extension
\[
\tag{43}\diagram{\begin{tikzcd}[ampersand replacement=\&,column sep=large]
\&\underline{\operatorname{End}}(M_0)\arrow[d,phantom,"\rotatebox{90}{$\simeq$}"]\&\&\underline{\operatorname{End}}(M_0)\arrow[d,phantom,"\rotatebox{90}{$\simeq$}"]\&\\
0\arrow[r]\&V_{\underline J}(M_0)\arrow[r]\&V_{\underline J}(M)\arrow[r]\&V_{\underline J}(M_0)\arrow[r]\&1
\end{tikzcd}}
\]
où pour $\mathcal A$ schéma \inserted{(ou faisceau)} en groupes comm. sur $S$, on pose
\[
\tag{44}V_{\underline J}(\mathcal A)\simeq\operatorname{Ker}(\mathcal A^{\underline J}\xrightarrow{\text{somme}}\mathcal A)
\]
(faisceau loc. isom. à $\mathcal A^2$), et $V_{\underline J}$ est un foncteur exact en $\mathcal A$, commutant bien sûr au chgt. de base. Pour $\mathcal A$ annulé par $2$, on a un isom
\[
\tag{45}V_{\underline J}(\mathcal A)\simeq\operatorname{Hom}(M_0,\mathcal A)\qquad\text{si }2\operatorname{id}_{\mathcal A}=0
\]
dû au fait que l'on a
\[
\tag{46}M_0\simeq\mathbb F_2^{\underline J}/\underset{\mathrm{diag}}{\underline{\mathbb F_2}}
\]
\unsure{On a|Donc à} une section \deleted{\unsure{la}} canonique partout non nulle de $V_J(M_0)$, correspondant à l'application d'inclusion $\underline J\unsure{\longleftrightarrow|=}M_0^*\hookrightarrow M_0$, ou encore à la section \unsure{évidente} $\operatorname{id}_{M_0}$ de $\underline{\operatorname{End}}(M_0)$, dont l'image inverse dans (43) est le torseur cherché.
\scanpage{68}
% archive foliation: 67
\deleted{Un épinglage de $M$ s'identifie donc à un}

On peut considérer, indépendamment de la question d'épingler $\underline J$, \deleted{celle de trouver} \inserted{l'épinglage} ``relatif'' \deleted{une section de} consistant en \deleted{une} la donnée d'une section du torseur $P_M$ sous $\underline{\operatorname{End}}(M_0)$, ce qui équivaut \inserted{donc} aussi à la donnée d'un morphisme
\[
\tag{47}\underline J\longrightarrow\deleted{\unsure{V_{\underline J}(}}M\deleted{\unsure{)}}
\]
qui relève le morphisme d'inclusion
\[
\underline J\hookrightarrow M_0
\]
et qui soit une section de $V_{\underline J}(M)$ i.e. qui soit à trace nulle (rel. à $S$).

On a en fait automatiquement un morphisme
\[
\underline J\longrightarrow M\setminus M_0
\]
donc en composant avec \deleted{\unsure{l'inclusion}}
\[
\diagram{\begin{tikzcd}[ampersand replacement=\&]
(\underline J\arrow[r,dotted]\& )M\arrow[r,hook]\&E\arrow[r]\&\Sigma_{\underline J}\simeq E/(1,\sigma_E)
\end{tikzcd}}
\]
on trouve un morphisme
\[
\tag{48}\underline J\longrightarrow\Delta_t
\]
et la condition de relèvement pour (47) signifie aussi que le composé de (48)
\scanpage{69}\mspage{13}
% archive foliation: 68
avec le morphisme canonique (étale de degré $2$)
\[
f_t:\Delta_t\longrightarrow\underline J
\]
soit l'identité -- i.e. (48) définit une \uline{face} du fibré en octaèdres combinatoires $\Delta_t$ sur $S$. Du coup l'octaèdre est $\underline J$-épinglé, en d'autres termes, il est can. isomorphe au \deleted{$\unsure{I}$} octaèdre standart défini par $\underline J$, \inserted{et une section de $\mu_4^*$} en prenant le revêtement biquadratique $\bar\Sigma_{\underline J,2}$ de $\bar\Sigma_{\underline J}$\ldots
\scanpage{70}\mspage{III}
% archive foliation: 69
Épinglage de Legendre

des courbes elliptiques
\scanpage{71}
% archive foliation: 70
\uline{Digression} Épinglages de revêtements par \deleted{\unsure{vecteurs} fibrés tangents} rev. kummériens des \unsure{Espaces|espaces} tangents (cas relatif)

\rule{\linewidth}{0.4pt}

Soient
\[
\tag{1}f:X\longrightarrow S
\]
un morphisme lisse de dim relative $1$, $t$ une section de $X$ sur $S$, $X^*=X\setminus t(S)$.

On considère des schémas \inserted{$X'$} finis sur $X$, lisses sur $S$ de dim relative $1$, munis d'une section $t'$ au dessus de $t$, et ``modérément ramifiés'' en $t'$ sur $X$, avec un degré de ramification $n$ -- ce qui signifie que loc pour la top. étale sur $X,X'$, $X'$ est \unsure{isom} au voisinage de tt point $x'$ de $t'(S)$ à un $X$-schéma Kummérien de $X$ au voisinage de $x=p(x')$, de la forme
\margin{\uline{NB} $n$ est premier aux car. résiduelles de $S$ !!!}
\margin{$p:X'\longrightarrow X$}
\[
X'_0=\operatorname{Spec}\underline O_X[z]/(z^n-\varphi)
\]
où $\varphi$ est une équation locale du diviseur relatif $t(S)\subset X$ au voisinage de $x$.

Sous ces conditions, via l'hom norme locale, on définit un isom. canonique
\[
\tag{2}T_{X'/\unsure{S},t'}^{\otimes n}\xrightarrow[\sim]{\kappa_{t'}}T_{X/S,t}
\]
\scanpage{72}
% archive foliation: 71

Si d'autre part $\underline G$ est un schéma en groupes fini étale sur $S$, et si $X'$ est muni d'une structure de $\underline G_X$-torseur au dessus de $X^*$, alors le stabilisateur \inserted{dans} $\underline G_{\deleted{\unsure{X}}}$ de la section $t'$ est un sous-schéma en groupes \inserted{$\underline G_{t'}$} de $\underline G$, \deleted{\unsure{schéma en groupes}} \inserted{qui} opère \inserted{fid.} sur \deleted{\unsure{comme si}} $T_{X'/S,t'}$ et dont l'image dans \deleted{$T_{X'/S,t'}$} $G_{mS}\simeq\underline{\operatorname{Aut}}_{\mathcal O_S}T_{X'/S,t'}$ est $\mu_{nS}$, d'où un monomorphisme canonique
\[
(3)\qquad \mu_{nS}\hookrightarrow\underline G.
\]
\margin{\uline{NB} On suppose \unsure{de même} que $\underline G$ opère sur $X'$, de façon que $X'^*=X'|X^*$ devienne un $\underline G_{X^*}$-torseur sur $X^*$.}

Par ailleurs, désignant par
\[
f^o:X^*\longrightarrow S
\]
le morphisme induit par $f$, le $\underline G_{X^*}$-torseur $X'^*$ définit une section
\[
(4)\qquad \unsure{\xi|\zeta}_{X'/S}\in H^0(S,R^1f^o_*f^{o*}(\underline G)).
\]
Tout ceci posé, on a :

\uline{Proposition} $(X',t')$ au dessus de $(X,t)$, avec l'opération de $\underline G$ dessus, est défini à isom unique près par la connaissance de (3)(4) et de (2). \deleted{\unsure{(ce dernier}} De façon précise,
\scanpage{73}
% archive foliation: 72

\deleted{pour une section}
considérons (pour $S,X,t$ fixés) la catégorie des $(X',t',\underline G)$, \deleted{où $X'$ est un $X$-schéma} \inserted{tels que dessus}, comme morphismes un
\[
(X',t',G)\mathbin{-}(X'_1,t'_1,G_1)
\]
\deleted{\unsure{étant}} couple \inserted{de} $X'\mathbin{-}X'_1$ ($X$-morphisme et de $G\mathbin{-}G_1$, compatibles dans un sens évident avec les sections marquées et les opérations.
\margin{\uline{NB} On n'exige pas que $n=n_1$. Pour $X',t',G$ fixés, $(X'_1,t'_1,G_1)$ est déterminé à isom can. près, quand on se donne le quotient $G\longrightarrow G_1$.}
Considérons d'autre part les \deleted{triples} systèmes
\[
(\underline G,\unsure{\deleted{\xi}}n,\deleted{\unsure{\xi}}i,\deleted{\unsure{G}}\unsure{\xi|\zeta},T,\kappa).
\]
où \inserted{$\underline G$ un schéma en groupes fini étale sur $S$,} $n\in\mathbb N^*$, $\unsure{\xi|\zeta}\in H^0(S,R^1f^o_*f^{o*}(\underline G))$ tels que les rev. locaux (ét. loc. sur $S$)) de $X$ qu'elle décrit soient modérément ramifiés d'ordre $n$, le long de $t$, $i:\mu_{nS}\hookrightarrow\underline G$ est un mono. de groupes, $T$ un Module inv. sur $S$, et $\kappa$ un isom
\[
T^{\otimes n}\xrightarrow{\sim}T_{X/S,t}
\]
avec la définition ad hoc des morphismes (\uline{NB} \unsure{tout hom.|tout un} $\underline G\mathbin{-}\underline G_1$ doit induire un épimorphisme $\mu_{nS}\mathbin{-}\mu_{n_1S}$, qui soit l'hom \deleted{\unsure{que}} $\lambda\longmapsto\lambda^{n/n_1}$ \ldots). On a un foncteur canonique de la première catégorie dans la
\scanpage{74}
% archive foliation: 73

seconde, et je dis que ce foncteur est une équivalence de catégories.

Le cas qui m'intéresse est celui où $\underline G=\mu_2\simeq\{\pm1\}_S$, et où $n=2$ --- donc où on s'intéresse aux rev. quadratiques $X'$ de $X/S$, étales au dessus de $X$ et ``totalement ramifiés le long de la section $t$'', avec $S$ à car. rés. $\ne2$. \deleted{\unsure{Ainsi}} Pour une section $\unsure{\xi|\zeta}\in H^0(S,R^1f^o_*(\mu_{2X^*}))$ fixée, satisfaisant à la condition de ramification totale, la catégorie des \deleted{\unsure{rev.}} revêtements \inserted{correspondants à cette $\unsure{\xi|\zeta}$} donc équivalente à celle des \deleted{\unsure{rev.}} couples $(T,\kappa)$ d'un Module inv. $T$ sur $S$, et d'un isom
\[
T^{\otimes2}\simeq T_{X,t}.
\]

Le cas qui nous intéresse pour les courbes elliptiques est celui où $X$ est de la forme $\Sigma_J^*$ (fibré en droites projectives associé à un rev. étale cubique $J$ de $S$), \inserted{déduit d'une courbe elliptique $X'$ comme $X'/\{\pm1\}$,} et où $t$ est la section de $\Sigma_J^{\unsure{*}}/S$ \deleted{\unsure{qui est l'}} image de la section nulle de $X'=E$ \unsure{\ldots}
\scanpage{75}
\mspage{1}
% archive foliation: 74

Soit $E$ une courbe elliptique --- on travaille sur un corps alg. clos $k$ pour simplifier le propos, mais la seule hypothèse essentielle est que la base soit à car. résiduelles $\ne2$. On désigne par \deleted{$\unsure{\Sigma|E}$}
\[
(1)\qquad J={}_2E(k)^*
\]
l'ens des pts de $E$ d'ordre $2$, de sorte que l'on a
\[
(2)\qquad E/\{\pm\operatorname{id}_E\}\simeq\Sigma_J
\qquad\text{\deleted{\unsure{sphère}} \inserted{droite projective} $J$-marquée.}
\]
\[
(3)\qquad\Sigma_J\simeq\mathbb P(V_J(k))
\qquad 0\mathbin{-}V_J(k)\mathbin{-}k^J\xrightarrow{\text{somme}}k\longrightarrow0
\]
La donnée de $E$ est équivalente à celle de $J$ i.e. de $\Sigma_J$ muni du diviseur étale $D_J$ de degré $3$, plus le point
\[
(4)\qquad t\in\Sigma_J^*=\Sigma_J\setminus D_J
\]
image de l'élément $0_E$, plus une droite $T$ \inserted{sur $k$} (savoir $T_{E,0}$), munie d'un isomorphisme
\[
(5)\qquad T^{\otimes2}\xrightarrow{\sim}T_{\Sigma,t}
\]
(\deleted{isomorphisme} \inserted{équivalence} de catégories-groupoïdes \ldots valable en fait sur base $S$ quelconque à car. rés. $\ne2$, et compatible aux changements de base \ldots).

1 Notons qu'on a un isomorphisme canonique
\[
(6)\qquad T_\Sigma\simeq\underline O_\Sigma(2)(\underline\omega)
\]
où $\underline\omega=\underline\omega(J)$ est l'ens. des deux ordres circulaires sur $J$, et où $L(\underline\omega)$ (pour un objet \unsure{$M$|$L$} \unsure{que|en} faisceaux abéliens etc) désigne $L\wedge_{\{\pm1\}}\underline\omega$.
\scanpage{76}
\mspage{2}
% archive foliation: 75

On a un candidat naturel (pour $\Sigma_{\deleted{J}}$ \inserted{$=\Sigma_J$}, \deleted{$D_J$} \inserted{$D=D_J$}, $t\in\Sigma^*$ donnés) comme droite $T$ munie de (5), savoir
\[
\deleted{(7)}\qquad T_0\simeq O_t(1)
\]
à cela près qu'il y a \uline{deux} isom. naturels
\[
\deleted{(8)}\qquad T_0^{\otimes2}=O_t(2)\xrightarrow{\sim}T_{\Sigma,t}\simeq O_t(2)(\underline\omega),
\]
correspondant aux deux éléments $\omega$ de $\underline\omega$, savoir
\[
(7)\qquad\deleted{\unsure{t}}\unsure{\xi|\zeta}\longmapsto\unsure{\xi|\zeta}\wedge\omega
\]
\margin{$T_0^{\otimes2}\xrightarrow{\sim}O_t(2)\simeq T_{\Sigma,t}\simeq\underline O_{\unsure{t}}(2)(\underline\omega)$}

On va lever cette ambiguïté, en remplaçant $T_0=O_t(1)$ par une variante tordue. La tordre par un torseur sous $\mu_2$ \inserted{$=\mu_2(k)$} ne nous avance pas, car \deleted{$\unsure{O}$} $T_0^{\otimes2}$ n'aura pas changé (à isom can. près). On va la tordre par un torseur (qu'on précisera) $Q_0$ sous $\mu_4=\mu_4(k)$ (qui peut donc s'interpréter comme l'ens des sommets d'un \uline{carré}, dont $\mu_4$ est le groupe des rotations, avec $\{i,-i\}=\mu_4^*$ comme \uuline{paire} de génér. privilégiés\ldots).

On posera donc
\[
(8)\qquad T_0=\underline O_t(1)\wedge_{\mu_4}Q_0
\]
d'où
\[
(9)\qquad T_0^{\otimes2}\simeq\underline O_t(2)\wedge_{\mu_2}D_{Q_0}
\]
où
\[
(10)\qquad D_{Q_0}=Q_0/\{\pm1\}
\]
est l'ens des deux diagonales du carré $Q_0$. Supposons donc qu'on se donne un carré \inserted{\unsure{combinatoire} (``sommets'')} $Q_{0S}$ avec des isom. \deleted{\unsure{$\operatorname{Aut}Q_0$} est le groupe des}
\[
(11)\qquad\left\{
\begin{array}{lll}
\underline\omega_{Q_0}&\text{(ens des orientations de $Q_0$)}&\simeq\mu_4^*=\{i,-i\}\\
D_{Q_0}&\text{(ens des diagonales de $Q_0$)}&\simeq\underline\omega
\end{array}\right.
\]
\scanpage{77}
% archive foliation: 76

Alors, définissant $T_0$ par (8) (pour $t\in\Sigma_J^*$ donné) on a un isom canonique
\[
(12)\qquad T_0^{\otimes2}\simeq T_{\Sigma,t}
\]
Ceci posé, on appellera \uline{épinglage de Legendre}
\deleted{\unsure{car il faut se rappeler que}}
\deleted{$\unsure{T_0}$ \unsure{ne correspond}}
de $E$, subordonné au carré $Q_0$ (muni des données (11)) \deleted{tout} \inserted{un} isomorphisme
\[
(13)\qquad T=T_{E,0}\xleftarrow{\substack{\kappa\\\sim}}T_0
\]
compatible avec les \unsure{hom. carrure}, i.e. rendant commutatif le diagramme
\diagram{$(14)\quad\begin{tikzcd}[ampersand replacement=\&,column sep=small,row sep=large]
T_0^{\otimes2}\arrow[rr,"{\substack{\kappa^{\otimes2}\\\sim}}"]\arrow[dr,"(12)"]\arrow[d,dashed,no head]\&\&T^{\otimes2}\arrow[dl,"(5)"]\\
{}\&T_{\Sigma,t}\&
\end{tikzcd}$}

Il y en a exactement deux, différant par le signe i.e. formant un torseur sous $\mu_2$.

Bien entendu, on veut pouvoir définir la notion d'épinglage de $E$ (en termes d'un $Q_0$ donné une fois pour toutes) pour $E$ variable --- mais alors $\underline\omega=\underline\omega_J$ n'est pas encore déterminé (puisque $J$ ne l'est pas). Donc la donnée d'un épinglage de Legendre sera définie comme \unsure{le couple|les couples} \inserted{$(\alpha,\kappa)$} formé d'un isom.
\[
(15)\qquad\alpha:\underline\omega_J\xleftarrow{\sim}D_{Q_0}=Q_0/\{\pm1\}
\]
\uline{et} d'un isom. $\kappa$ (13) satisfaisant la compatibilité (14).
